\begin{statementbox}
\begin{center}\begin{tikzpicture}[font=\small,>=stealth,line join=round]
\draw[gray!40,rounded corners=3pt] (0,0) rectangle (14.2,8.3);\node at (7.1,7.700000000000001) {三种写法，对应同一条抛物线};
\begin{scope}[shift={(1.7000000000000002,2.9642857142857144)},x=1.2cm,y=0.6571428571428571cm]\draw[->,gray] (-0.5,0)--(4.5,0) node[right]{$x$};\draw[->,gray] (0,-2)--(0,5) node[above]{$y$};\draw[blue,thick,solid,domain=-0.2:4.2,samples=180] plot(\x,{(\x-2)^2-1});\fill[orange!85!black] (0,3) circle(2.2pt);\fill[orange!85!black] (1,0) circle(2.2pt);\fill[orange!85!black] (3,0) circle(2.2pt);\fill[orange!85!black] (2,-1) circle(2.2pt);\draw[violet,dashed] (2,-.95)--(2,4.5);\node[left] at (0,3){$(0,3)$};\node[below=5pt] at (1,0){$1$};\node[below=5pt] at (3,0){$3$};\node[below=4pt] at (2,-1){$(2,-1)$};\end{scope}\node[anchor=west,align=left,text width=5.0cm] at (8.4,4.0) {$x^2-4x+3$\\[9pt]$=(x-2)^2-1$\\[9pt]$=(x-1)(x-3)$\\[9pt]顶点 $(2,-1)$\\[9pt]交点 $(1,0),(3,0)$};
\node[text width=12.6cm,align=center] at (7.1,.65) {一般式读纵轴交点，顶点式读顶点，两根式读横轴交点。};
\end{tikzpicture}\end{center}


\textbf{\textcolor{CStatement}{这条结论解决什么问题？}}
同一条抛物线，题目可能告诉我们三个系数、顶点或与 $x$ 轴的交点。
如果每次都从 $ax^2+bx+c$ 开始列方程，就容易绕远路。
\textbf{已知什么，就先写最容易放进这些已知条件的式子。}

\medskip
\end{statementbox}\clearpage\begin{statementbox}\textbf{\textcolor{CStatement}{成立条件}}
二次函数必须有 $a\ne0$，也就是 $x^2$ 前面的系数不能是 $0$；以下讨论的 $a,b,c,x$ 均为实数。

\medskip
\textbf{\textcolor{CStatement}{三种写法，一条抛物线}}
\begin{description}[style=nextline,leftmargin=2.8em,labelsep=0.5em,itemsep=0.55em]
\item[\textbf{（一） 一般式：容易看系数和纵轴交点}]
\[
\boxed{f(x)=ax^2+bx+c\quad(a\ne0).}
\]
代入 $x=0$，就得 $f(0)=c$，所以图象经过 $(0,c)$。当题目给出多个点的坐标时，这种写法通常方便列方程。
\item[\textbf{（二） 顶点式：一眼看出最高点或最低点}]
\[
\boxed{f(x)=a(x-h)^2+k,\qquad
h=-\frac{b}{2a},\quad k=c-\frac{b^2}{4a}.}
\]
顶点是 $(h,k)$，对称轴是 $x=h$。$a>0$ 时开口向上，最小值是 $k$；$a<0$ 时开口向下，最大值是 $k$。
顶点式\textbf{对每一个二次函数都能写出来}。
\item[\textbf{（三） 两根式（也叫交点式）：一眼看出与 $x$ 轴的交点}]
如果 $f(x)=0$ 有实数解 $x_1,x_2$（允许两个解相同），那么
\[
\boxed{f(x)=a(x-x_1)(x-x_2).}
\]
图象与 $x$ 轴的交点为 $(x_1,0),(x_2,0)$；如果 $x_1=x_2$，它只在这一点碰到 $x$ 轴。
\end{description}

\medskip
\end{statementbox}\clearpage\begin{statementbox}\textbf{\textcolor{CStatement}{什么时候能用两根式？}}
\begin{center}\begin{tikzpicture}[font=\small,>=stealth,line join=round]
\draw[gray!40,rounded corners=3pt] (0,0) rectangle (14.2,8.5);\node at (7.1,7.9) {开口向上：顶点高低决定实数交点个数};
\node at (2.4,6.9) {$\Delta=4>0$};\begin{scope}[shift={(2.4,3.8)},x=0.85cm,y=0.6cm]\draw[->,gray] (-2,0)--(2,0) node[right]{$x$};\draw[->,gray] (0,-2)--(0,4) node[above]{$y$};\draw[blue,thick,solid,domain=-1.7:1.7,samples=160] plot(\x,{(\x)^2-1});\fill[orange!85!black] (-1,0) circle(2pt);\fill[orange!85!black] (1,0) circle(2pt);\fill[orange!85!black] (0,-1) circle(2pt);\end{scope}\node[align=center,text width=4cm] at (2.4,1.7) {$y=x^2-1$\\两个交点};\node at (7.050000000000001,6.9) {$\Delta=0$};\begin{scope}[shift={(7.050000000000001,3.8)},x=0.85cm,y=0.6cm]\draw[->,gray] (-2,0)--(2,0) node[right]{$x$};\draw[->,gray] (0,-2)--(0,4) node[above]{$y$};\draw[teal,thick,solid,domain=-1.7:1.7,samples=160] plot(\x,{(\x)^2});\fill[orange!85!black] (0,0) circle(2pt);\end{scope}\node[align=center,text width=4cm] at (7.050000000000001,1.7) {$y=x^2$\\一个交点};\node at (11.7,6.9) {$\Delta=-4<0$};\begin{scope}[shift={(11.7,3.8)},x=0.85cm,y=0.6cm]\draw[->,gray] (-2,0)--(2,0) node[right]{$x$};\draw[->,gray] (0,-2)--(0,4) node[above]{$y$};\draw[violet,thick,solid,domain=-1.7:1.7,samples=160] plot(\x,{(\x)^2+1});\fill[orange!85!black] (0,1) circle(2pt);\end{scope}\node[align=center,text width=4cm] at (11.7,1.7) {$y=x^2+1$\\没有交点};
\node[text width=12.6cm,align=center] at (7.1,.65) {同尺度，$a=1$：$k=-1,0,1$，恒有$\Delta=-4ak$。};
\end{tikzpicture}\end{center}


令 $\Delta=b^2-4ac$，则
\[
\begin{array}{c|c|c}
\Delta>0 & \text{两个不同的实数解} & \text{与 }x\text{ 轴有两个交点}\\
\Delta=0 & \text{两个相同的实数解} & \text{与 }x\text{ 轴只有一个交点}\\
\Delta<0 & \text{没有实数解} & \text{与 }x\text{ 轴没有交点}
\end{array}
\]
只有 $\Delta\ge0$，才可在\textbf{实数范围内}写出两根式。

\medskip
\textbf{\textcolor{CStatement}{三个值得连在一起记的关系}}
\[
\boxed{\;x_1+x_2=-\frac ba=2h,\qquad
x_1x_2=\frac ca,\qquad
\Delta=-4ak.\;}
\]
前两式需要有实数解才能把 $x_1,x_2$ 理解为图上的交点；最后一式对每个二次函数都成立。
\textbf{特别是：有两个交点时，对称轴一定经过它们的中点。}
\end{statementbox}
